% JUMP TO LINE 60, 75
\documentclass[preview, margin=0.6in]{standalone}
\usepackage[letterpaper,portrait,top=0.4in, left=0.6in, right=0.6in, bottom=1in]{geometry}

\usepackage{amsmath, amsfonts, amsthm, amssymb}
\usepackage{graphicx, float}
\usepackage{mathtools}
\usepackage{titlesec}
\usepackage{interval}
\usepackage{hyperref}
\usepackage{siunitx}
\usepackage{titling}
\usepackage{vwcol}
\usepackage{setspace}
\usepackage{empheq}
\usepackage{cancel}
\usepackage{esdiff}
\usepackage{multicol}
\usepackage{mdframed}
\usepackage{esdiff}
\usepackage{tikzsymbols}
\usepackage{multicol}
\usepackage{tikz}
\usepackage{varwidth}
\usepackage{parskip}
\usepackage{pgfplots}
\usepackage{mlmodern}
\pgfplotsset{compat=1.18}
\intervalconfig {
	soft open fences
}

\newcommand{\alignedintertext}[1]{%
  \noalign{%
    \vtop{\hsize=\linewidth#1\par
    \expandafter}%
    \expandafter\prevdepth\the\prevdepth
  }%
}

\newtheorem{lemma}{Lemma}

\newcommand*{\problem}[1]{\section*{Problem #1}}
\newcommand*{\aps}{\section*{AP Corner}}
\newcommand*{\deriv}[1][x]{\ensuremath{\dfrac{\mathrm{d}}{\mathrm{d}#1}}}
\newcommand*{\floor}[1]{\ensuremath{\lfloor #1\rfloor}}
\newcommand*{\lheqzero}{\ensuremath{\underset{\text{L'H}}{\overset{\left[\frac00\right]}{=}}}}
\newcommand*{\lheqinfty}{\ensuremath{\underset{\text{L'H}}{\overset{\left[\frac{\infty}{\infty}\right]}{=}}}}

\DeclareMathOperator{\DNE}{DNE}
\DeclareMathOperator{\sgn}{sgn}

\DeclareMathOperator{\arccsc}{arccsc}
\DeclareMathOperator{\arcsec}{arcsec}
\DeclareMathOperator{\arccot}{arccot}

%opening

\title{\vspace*{-40pt}Classwork Set A}
\author{Jayden Li}
\date{\today}
% \allowdisplaybreaks
\postdisplaypenalty=100000

\begin{document}
\setstretch{1.25}
\fontsize{12pt}{12pt}\selectfont
\setlength{\abovedisplayskip}{\abovedisplayskip/2}
\setlength{\belowdisplayskip}{\belowdisplayskip/2}
\setlength{\parindent}{0pt}
\setlength{\parskip}{2ex plus 0.5ex minus 0.2ex}
\maketitle

\problem{2}

\begin{itemize}
\item[(b)]
	Let us fix the vector $v_{3}=(r,0)$. We need to show that $v_i+v_j+v_{3}=0$ if and only if the angle between $v_i$ and $v_j$ is $120^\circ=2\pi/3$, and the angle between $v_3$ and $v_i$ is also $120^\circ=2\pi/3$.

	\begin{proof}
		\begin{align*}
			&v_i+v_j+v_{3}=0 \\
			\iff{}& v_i+v_j=-v_{3}= \begin{bmatrix}
				-r \\ 0
			\end{bmatrix}
			\intertext{Let $\theta\in [-\pi,\pi]$ be the angle between $v_3$ and $v_i$, and let $\varphi\in[-\pi,\pi]$ be the angle between $v_3$ and $v_j$. The radius of the clock is $r$. We can find the components of the two vectors: $v_i=(r\cos\theta,r\sin\theta)$ and $v_j=(r\cos\varphi,r\sin\varphi)$.}
			\iff{}&
			\begin{bmatrix}
				r\cos\theta \\ r\sin\theta
			\end{bmatrix}
			+\begin{bmatrix}
				r\cos\varphi \\ r\sin\varphi
			\end{bmatrix}
			=\begin{bmatrix}
				-r \\ 0
			\end{bmatrix} \\
			\iff{}&
			\left\{\begin{aligned}
				r\cos\theta+r\cos\varphi&=-r \\
				r\sin\theta+r\sin\varphi&=0 \\
			\end{aligned}\right.
			\intertext{Notice: $r\sin\theta+r\sin\varphi=0\iff \sin\theta=-\sin\varphi$. Since $\theta,\varphi\in[-\pi,\pi]$, we have $\varphi=-\theta$.}
			\iff{}& r\cos\theta+r\cos(-\theta)=-r \\
			\iff{}& 2\cos\theta=-1 \\
			\iff{}& \cos\theta=-\frac12 \\
			\iff{}& \theta=\pm\frac{2\pi}{3}
		\end{align*}
		Therefore, the angle between $v_3$ and $v_i$ is $\theta=\pm2\pi/3$, which can be either $2\pi/3$ or $-2\pi/3$. Either way, $\varphi=-\theta=\mp2\pi/3$. As such, \textbf{the angle between $v_i$ and $v_3$ is $120^\circ$, as is the angle between $v_j$ and $v_3$}. The angle between $v_i$ and $v_j$ is $2\pi-2\pi/3-2\pi/3=2\pi/3=120^\circ$. 
	\end{proof}
\end{itemize}

\problem{3}
\begin{itemize}
	\item[(a)] True. If $\theta$ is the direction of $u,v$ and $m$ is the magnitude then $u=v=(m\cos\theta,m\sin\theta)$.
	\item[(b)] False. If $a$ is negative, then $\sqrt{2}a$ is negative but magnitude $\lVert (a,b) \rVert $ must be positive.
	\item[(c)] True. $u$ and $v$ point in the same direction. $u$ can be defined as $u=v/\lVert v \rVert $, and therefore:
		\begin{equation*}
		    \lVert v \rVert u
			=\cancel{\lVert v \rVert}\cdot \frac{v}{\cancel{\lVert v \rVert}}
			=v
		\end{equation*}
	\item[(d)] True. Since $\hat i=(1,0)$ and $\hat j=(0,1)$ are the standard unit vectors:
		\begin{equation*}
		    \lVert u \rVert 
			=\lVert a\hat i+b\hat j \rVert 
			=\lVert (a,0)+(0,b) \rVert 
			=\lVert (a,b) \rVert 
			=\sqrt{a^2+b^2}
			=1
			\implies a^2+b^2=1
		\end{equation*}
	\item[(e)] True. If $v=a\hat i+b\hat j=0$, then $v=(0,0)\iff a=0,b=0 \implies a=-b$.
\end{itemize}

\problem{4}
9. Let $c=-1$ and $\mathbf{v}=(1,0)$. $\lVert c \mathbf{v} \rVert =\lVert (-1,0) \rVert =1 \neq -1 =(-1) \lVert (1,0) \rVert =c \lVert \mathbf{v} \rVert $.
\end{document}
